%%%%%%%%%%%%%%%%%%%%%%%%%%%%%%%%%%%%%%%%%%%%%%%%%%%%%%%%%%%%%%%%%%%%%%%%%%%%%%%%
\begin{frame}[fragile]\frametitle{}
\begin{center}
{\Large State, Memory and Runtime Control}
\end{center}
\end{frame}
%%%%%%%%%%%%%%%%%%%%%%%%%%%%%%%%%%%%%%%%%%%%%%%%%%%%%%%%%%%%%%%%%%%%%%%%%%%%%%%%
\begin{frame}[fragile]\frametitle{Is Each Node an Agent or Is the Graph an Agent?}
      \begin{itemize}
        \item \textbf{The Graph is the Agent} - not individual nodes
        \item \textbf{Nodes are}: Individual processing units/functions
        \begin{itemize}
            \item Can be simple Python functions
            \item Can be LLM calls
            \item Can be tool invocations
            \item Can be sub-agents themselves
        \end{itemize}
        \item \textbf{The Graph as Agent}:
        \begin{itemize}
            \item The complete workflow represents the agent's behavior
            \item State flows through nodes, creating agent "memory"
            \item Control flow (edges) represents agent's decision-making
            \item Overall graph exhibits autonomous, goal-directed behavior
        \end{itemize}
        \item \textbf{Analogy}: Nodes are like neurons, the graph is the brain
        \item Individual nodes are stateless functions; the graph maintains state
      \end{itemize}
\end{frame}
%%%%%%%%%%%%%%%%%%%%%%%%%%%%%%%%%%%%%%%%%%%%%%%%%%%%%%%%%%%%%%%%%%%%%%%%%%%%%%%%
\begin{frame}[fragile]\frametitle{What Can Nodes Do? - Part 1}
      \begin{itemize}
        \item \textbf{Nodes are Extremely Flexible} - any Python function
        \item \textbf{1. Simple LLM Calls}:
        \begin{itemize}
            \item Generate text, classify, extract information
            \item Example: \texttt{llm.invoke("Summarize this text")}
        \end{itemize}
        \item \textbf{2. RAG Operations}:
        \begin{itemize}
            \item Retrieve documents from vector database
            \item Rerank results
            \item Generate answers with retrieved context
            \item Example: Vector search \(\rightarrow\) Context \(\rightarrow\) LLM generation
        \end{itemize}
        \item \textbf{3. API Calls}:
        \begin{itemize}
            \item External service calls (weather, database, CRM)
            \item Example: Fetch customer data from Salesforce API
        \end{itemize}
      \end{itemize}
\end{frame}
%%%%%%%%%%%%%%%%%%%%%%%%%%%%%%%%%%%%%%%%%%%%%%%%%%%%%%%%%%%%%%%%%%%%%%%%%%%%%%%%
\begin{frame}[fragile]\frametitle{What Can Nodes Do? - Part 2}
      \begin{itemize}
        \item \textbf{4. Tool-Calling Agents}:
        \begin{itemize}
            \item Node can itself be an agent with tools
            \item Example: Research node that uses search, calculator, Wikipedia tools
            \item Inner agent makes tool decisions, outer graph manages workflow
        \end{itemize}
        \item \textbf{5. Multi-Agent Nodes}:
        \begin{itemize}
            \item Node can coordinate multiple sub-agents
            \item Example: "Code Review" node with architect, tester, security agents
        \end{itemize}
        \item \textbf{6. Any Computation}:
        \begin{itemize}
            \item Data processing, calculations, file operations
            \item Validation, formatting, business logic
        \end{itemize}
        \item \textbf{Key Point}: Nodes are building blocks - combine them creatively
      \end{itemize}
\end{frame}
%%%%%%%%%%%%%%%%%%%%%%%%%%%%%%%%%%%%%%%%%%%%%%%%%%%%%%%%%%%%%%%%%%%%%%%%%%%%%%%%
\begin{frame}[fragile]\frametitle{Custom Node Types}

Any callable is a node, so a class with \texttt{\_\_call\_\_} can carry its own configuration (model, prompt, tools).

      \begin{lstlisting}[language=Python]
from langchain_groq import ChatGroq
from langgraph.graph import StateGraph, MessagesState, START

class MyCustomNode:
    def __init__(self, llm):
        self.llm = llm

    def __call__(self, state: MessagesState):
        # Access the state, do the work, return only what changed
        response = self.llm.invoke(state["messages"])
        return {"messages": [response]}

llm = ChatGroq(model="openai/gpt-oss-20b")
graph_builder = StateGraph(MessagesState)
graph_builder.add_node("custom_node", MyCustomNode(llm))
graph_builder.add_edge(START, "custom_node")
      \end{lstlisting}
\end{frame}
%%%%%%%%%%%%%%%%%%%%%%%%%%%%%%%%%%%%%%%%%%%%%%%%%%%%%%%%%%%%%%%%%%%%%%%%%%%%%%%%
\begin{frame}[fragile]\frametitle{Advanced Conditional Edges}

A routing function reads the state and returns the name of the next route; the path map turns routes into nodes.

% [code-font: forced \scriptsize | overflow at the default size]
      \begin{lstlisting}[language=Python, basicstyle=\scriptsize\ttfamily]
from typing import Annotated, Literal
from langchain_core.messages import AnyMessage
from langgraph.graph import END
from langgraph.graph.message import add_messages
from pydantic import BaseModel, Field

class AgentState(BaseModel):
    messages: Annotated[list[AnyMessage], add_messages] = Field(default_factory=list)
    status: str = "RUNNING"
    error_count: int = 0

def route_by_status(state: AgentState) -> Literal["process", "retry", "error", "end"]:
    if state.status == "SUCCESS":
        return "end"
    if state.status == "ERROR":
        return "error" if state.error_count >= 3 else "retry"
    return "process"

workflow.add_conditional_edges(
    "check_status", route_by_status,
    {"process": "execute_tool", "retry": "retry_handler",
     "error": "error_handler", "end": END})
      \end{lstlisting}
\end{frame}
%%%%%%%%%%%%%%%%%%%%%%%%%%%%%%%%%%%%%%%%%%%%%%%%%%%%%%%%%%%%%%%%%%%%%%%%%%%%%%%%
\begin{frame}[fragile]\frametitle{State Updates: Overwrite vs Reducers}
      \begin{itemize}
        \item By default, a value returned by a node \textbf{overwrites} the field in the state.
        \item That silently drops earlier data: returning \texttt{[4]} for a list holding \texttt{[1, 2, 3]} leaves \texttt{[4]}.
        \item A \textbf{reducer function} declares how an update is combined with the current value.
        \item Attach it with \texttt{Annotated[type, reducer]}; LangGraph then calls \texttt{reducer(current, update)}.
        \item Built-ins: \texttt{add\_messages} for chat history, \texttt{operator.add} for lists and numbers.
        \item Reducers also make \textbf{parallel} branches safe: each branch's update is merged, not lost.
      \end{itemize}
\end{frame}
%%%%%%%%%%%%%%%%%%%%%%%%%%%%%%%%%%%%%%%%%%%%%%%%%%%%%%%%%%%%%%%%%%%%%%%%%%%%%%%%
\begin{frame}[fragile]\frametitle{Reducers: The Overwrite Problem}

By default, the dictionary returned by a node overwrites the original data, which may lead to unexpected results.

\begin{lstlisting}[language=Python]
from typing_extensions import TypedDict

class MyState(TypedDict):
    messages: list                      # no reducer: overwrite

def fn1(state: MyState):
    return {"messages": [4]}

# graph with a single node fn1
r = graph.invoke({"messages": [1, 2, 3]})
# Result: {"messages": [4]}   instead of [1, 2, 3, 4]
\end{lstlisting}
\end{frame}
%%%%%%%%%%%%%%%%%%%%%%%%%%%%%%%%%%%%%%%%%%%%%%%%%%%%%%%%%%%%%%%%%%%%%%%%%%%%%%%%
\begin{frame}[fragile]\frametitle{Reducers: The Fix}

Either read and extend the old value by hand, or declare a reducer once on the field.

\begin{lstlisting}[language=Python]
# Manual: fetch the old value and return old + new
def fn1(state: MyState):
    old = state.get("messages", [])
    return {"messages": old + [4]}

# Declarative: a reducer on the field (a custom one, or operator.add)
from typing import Annotated

def concat_lists(original: list, new: list) -> list:
    return original + new

class MyState(TypedDict):
    messages: Annotated[list, concat_lists]

def fn1(state: MyState):
    return {"messages": [4]}            # only the new item

r = graph.invoke({"messages": [1, 2, 3]})
# Result: {"messages": [1, 2, 3, 4]}
\end{lstlisting}
\end{frame}
%%%%%%%%%%%%%%%%%%%%%%%%%%%%%%%%%%%%%%%%%%%%%%%%%%%%%%%%%%%%%%%%%%%%%%%%%%%%%%%%
\begin{frame}[fragile]\frametitle{Checkpointing: Saving the State at Every Step}
      \begin{itemize}
        \item A \textbf{checkpointer} saves a snapshot of the state after each step, keyed by \texttt{thread\_id}.
        \item \textbf{Persistence}: stop and resume long workflows across sessions.
        \item \textbf{Resilience}: recover from failures without losing progress.
        \item \textbf{Time travel}: inspect or replay any earlier step.
        \item \textbf{Human-in-the-loop}: pause, wait, resume.
        \item \textbf{Backends}: in-memory for development, SQLite for local files, Postgres or Redis for production.
      \end{itemize}
% [code-font: forced \scriptsize | overflow at the default size]
\begin{lstlisting}[language=Python, basicstyle=\scriptsize\ttfamily]
import sqlite3
from langgraph.checkpoint.memory import MemorySaver
from langgraph.checkpoint.sqlite import SqliteSaver   # pip install langgraph-checkpoint-sqlite

memory = MemorySaver()                                 # development only
# memory = SqliteSaver(sqlite3.connect("agent_state.db", check_same_thread=False))

app = graph_builder.compile(checkpointer=memory)
config = {"configurable": {"thread_id": "user-123"}}
app.invoke({"messages": [("human", "Hi!")]}, config)   # state saved automatically
\end{lstlisting}
\end{frame}
%%%%%%%%%%%%%%%%%%%%%%%%%%%%%%%%%%%%%%%%%%%%%%%%%%%%%%%%%%%%%%%%%%%%%%%%%%%%%%%%
\begin{frame}[fragile]\frametitle{Interacting with Checkpoints}
      \begin{itemize}
        \item With a checkpointer you can inspect, modify and travel through any thread's history.
        \item A thread is identified by \texttt{thread\_id}; each snapshot by its \texttt{checkpoint\_id}.
      \end{itemize}
\begin{lstlisting}[language=Python]
config = {"configurable": {"thread_id": "user-123"}}

# Latest state of the conversation
snapshot = app.get_state(config)
print(snapshot.values)        # the state fields
print(snapshot.next)          # nodes still to run, () when finished

# Every snapshot, newest first (an iterator: wrap it in list)
history = list(app.get_state_history(config))
print(f"History has {len(history)} steps.")

# Write into the state directly, e.g. to inject a correction
app.update_state(config, {"messages": [("user", "Actually, a different question.")]})
\end{lstlisting}
\end{frame}
%%%%%%%%%%%%%%%%%%%%%%%%%%%%%%%%%%%%%%%%%%%%%%%%%%%%%%%%%%%%%%%%%%%%%%%%%%%%%%%%
\begin{frame}[fragile]\frametitle{Time Travel: Replay or Fork from a Checkpoint}
      \begin{itemize}
        \item Each snapshot carries a config with its \texttt{checkpoint\_id}.
        \item Invoke with \texttt{None} and that config to replay from that point.
        \item \texttt{update\_state} on an old snapshot forks a new branch; the original history is kept.
      \end{itemize}
\begin{lstlisting}[language=Python]
history = list(app.get_state_history(config))

# Pick the snapshot just before a given node ran
before_send = next(s for s in history if s.next == ("send_email",))

# Replay from there (nodes after that point run again)
app.invoke(None, before_send.config)

# Or fork: edit the old state, then continue on a new branch
forked = app.update_state(before_send.config, {"draft": "A better draft"})
app.invoke(None, forked)
\end{lstlisting}
\end{frame}
%%%%%%%%%%%%%%%%%%%%%%%%%%%%%%%%%%%%%%%%%%%%%%%%%%%%%%%%%%%%%%%%%%%%%%%%%%%%%%%%
\begin{frame}[fragile]\frametitle{Message Management in LangGraph}
      \begin{itemize}
        \item Message history grows with every turn: cost, latency and context limits follow.
        \item \texttt{RemoveMessage(id=...)}, returned from a node, deletes a message through the \texttt{add\_messages} reducer.
        \item \texttt{trim\_messages} keeps the last N messages or tokens.
        \item Summarise old turns into one message, then remove the originals.
        \item Returning a shorter list does not work: the reducer appends, it does not replace.
      \end{itemize}
\end{frame}
%%%%%%%%%%%%%%%%%%%%%%%%%%%%%%%%%%%%%%%%%%%%%%%%%%%%%%%%%%%%%%%%%%%%%%%%%%%%%%%%
\begin{frame}[fragile]\frametitle{Message Deletion and Trimming}

\begin{lstlisting}[language=Python]
from langchain_core.messages import RemoveMessage, ToolMessage, trim_messages

def drop_tool_messages(state):
    # Delete processed tool outputs by id
    return {"messages": [RemoveMessage(id=m.id)
                         for m in state["messages"]
                         if isinstance(m, ToolMessage)]}

def keep_latest(state, max_messages=50):
    # Remove everything except the last max_messages
    old = state["messages"][:-max_messages]
    return {"messages": [RemoveMessage(id=m.id) for m in old]}

# Trim only what is sent to the model, keep full history in the state
def call_model(state):
    recent = trim_messages(state["messages"], max_tokens=20,
                           token_counter=len, strategy="last")
    return {"messages": [llm.invoke(recent)]}
\end{lstlisting}
\end{frame}
%%%%%%%%%%%%%%%%%%%%%%%%%%%%%%%%%%%%%%%%%%%%%%%%%%%%%%%%%%%%%%%%%%%%%%%%%%%%%%%%
\begin{frame}[fragile]\frametitle{Error Handling Patterns}
      \begin{itemize}
        \item \textbf{Robust agents must handle unexpected failures}, like API errors or invalid LLM outputs.
        \item \textbf{Pattern 1: Try/Except in Nodes}: wrap node logic and return an error message in the state.
        \item \textbf{Pattern 2: Fallback Edges}: a router sends an 'ERROR' status to a recovery or reporting node.
        \item \textbf{Pattern 3: Retry Policy}: \texttt{add\_node(..., retry\_policy=RetryPolicy(max\_attempts=3))} retries transient errors with backoff.
        \item \textbf{State Tracking}: count errors in the state; after N retries route to a final failure node.
      \end{itemize}
% [code-font: forced \tiny | overflow at the default size]
% [code-font: blank lines removed from the listing to fit the frame]
\begin{lstlisting}[language=Python, basicstyle=\tiny\ttfamily]
def tool_node(state):
    try:
        result = call_external_api(state)
        return {"result": result, "status": "SUCCESS"}
    except Exception as e:
        return {"error_message": str(e), "status": "ERROR"}
def router(state):
    return "error_handler_node" if state.get("status") == "ERROR" else "next_step_node"
workflow.add_conditional_edges("tool_node", router)
\end{lstlisting}
\end{frame}
%%%%%%%%%%%%%%%%%%%%%%%%%%%%%%%%%%%%%%%%%%%%%%%%%%%%%%%%%%%%%%%%%%%%%%%%%%%%%%%%
\begin{frame}[fragile]\frametitle{Testing LangGraph Applications}
      \begin{itemize}
        \item \textbf{Unit Testing Nodes}:
        \begin{itemize}
            \item Nodes are just Python functions. You can test them in isolation by creating a sample input state and asserting the output dictionary is correct.
            \item This allows you to verify the logic of each step independently.
        \end{itemize}
        \item \textbf{Integration Testing the Graph}:
        \begin{itemize}
            \item Test the compiled graph (\texttt{app}) by invoking it with a representative input.
            \item Assert that the final state is what you expect. This tests the interaction between nodes and the correctness of the routing logic.
            \item Use \texttt{app.stream()} to check the state at intermediate steps.
        \end{itemize}
        \item \textbf{Mocking}: Use mocking libraries (like \texttt{unittest.mock}) to replace external dependencies like LLMs or APIs during tests. This makes tests faster, cheaper, and more predictable.
      \end{itemize}
\end{frame}
%%%%%%%%%%%%%%%%%%%%%%%%%%%%%%%%%%%%%%%%%%%%%%%%%%%%%%%%%%%%%%%%%%%%%%%%%%%%%%%%
\begin{frame}[fragile]\frametitle{Parallel Execution: Fan-Out and Fan-In}
      \begin{itemize}
        \item \textbf{What}: nodes that start from the same node run in the same step, in parallel.
        \item \textbf{Use cases}: concurrent API calls, several data sources, several variations of an answer.
        \item \textbf{Needs a reducer}: parallel branches write the same key, so declare how updates merge.
      \end{itemize}
% [code-font: forced \scriptsize | overflow at the default size]
\begin{lstlisting}[language=Python, basicstyle=\scriptsize\ttfamily]
import operator
from typing import Annotated
from typing_extensions import TypedDict

class SearchState(TypedDict):
    query: str
    results: Annotated[list[str], operator.add]   # merge parallel writes

builder = StateGraph(SearchState)
builder.add_node("search_web", search_web)        # each returns {"results": [...]}
builder.add_node("search_db", search_db)
builder.add_node("search_docs", search_docs)
builder.add_node("aggregate", aggregate_results)

for n in ("search_web", "search_db", "search_docs"):
    builder.add_edge(START, n)                    # fan-out
builder.add_edge(["search_web", "search_db", "search_docs"], "aggregate")  # fan-in
\end{lstlisting}
\end{frame}
%%%%%%%%%%%%%%%%%%%%%%%%%%%%%%%%%%%%%%%%%%%%%%%%%%%%%%%%%%%%%%%%%%%%%%%%%%%%%%%%
\begin{frame}[fragile]\frametitle{Parallel Node Execution: Rules of the Road}
      \begin{itemize}
        \item END ends one route, not the whole graph: the run finishes when no node is left to run.
        \item Nodes in the same step run in parallel, and their completion order is not guaranteed.
        \item A list in \texttt{add\_edge([a, b], "merge")} makes \texttt{merge} wait for all of them.
        \item Never let parallel nodes depend on each other's output in the same step.
        \item Keep merged fields order-independent, or sort them in the aggregate node.
      \end{itemize}
\begin{lstlisting}[language=Python]
graph.add_edge(["left1", "right3"], "merge")   # join: waits for both branches
\end{lstlisting}
\end{frame}
%%%%%%%%%%%%%%%%%%%%%%%%%%%%%%%%%%%%%%%%%%%%%%%%%%%%%%%%%%%%%%%%%%%%%%%%%%%%%%%%
\begin{frame}[fragile]\frametitle{Parallel Tool Execution Inside a Node}

Concurrency can also live inside one node, for example several tools called with \texttt{asyncio}.

      \begin{lstlisting}[language=Python]
import asyncio

async def parallel_tools_execution(state: AgentState) -> dict:
    """Parallel execution of multiple tools"""
    tools = identify_required_tools(state.current_input)

    async def run_tool(tool):
        return {tool.name: await tool.ainvoke(state.current_input)}

    # Execute all tools in parallel, then merge the results
    results = await asyncio.gather(*[run_tool(t) for t in tools])
    tools_output = {}
    for r in results:
        tools_output.update(r)

    return {"tools_output": tools_output, "status": "SUCCESS"}
      \end{lstlisting}
\end{frame}
%%%%%%%%%%%%%%%%%%%%%%%%%%%%%%%%%%%%%%%%%%%%%%%%%%%%%%%%%%%%%%%%%%%%%%%%%%%%%%%%
\begin{frame}[fragile]\frametitle{Dynamic Work: Map-Reduce with Send}
      \begin{itemize}
        \item The graph structure is fixed at compile time; what can vary is \textbf{how many} branches run.
        \item A routing function returns a list of \texttt{Send(node, payload)}: one parallel run per item.
        \item \textbf{Map}: split the work. \textbf{Process}: run the same node on each piece. \textbf{Reduce}: a reducer collects the results.
      \end{itemize}
% [code-font: forced \scriptsize | overflow at the default size]
\begin{lstlisting}[language=Python, basicstyle=\scriptsize\ttfamily]
from langgraph.types import Send

class MRState(TypedDict):
    docs: list[str]
    summaries: Annotated[list[str], operator.add]   # reduce

def fan_out(state: MRState):
    return [Send("summarize", {"doc": d}) for d in state["docs"]]   # map

def summarize(payload: dict):
    return {"summaries": [llm.invoke(f"Summarise: {payload['doc']}").content]}

g = StateGraph(MRState)
g.add_node("summarize", summarize)
g.add_conditional_edges(START, fan_out, ["summarize"])
g.add_edge("summarize", END)
\end{lstlisting}
\end{frame}
%%%%%%%%%%%%%%%%%%%%%%%%%%%%%%%%%%%%%%%%%%%%%%%%%%%%%%%%%%%%%%%%%%%%%%%%%%%%%%%%
\begin{frame}[fragile]\frametitle{Dynamic Routing: Command(update, goto)}
      \begin{itemize}
        \item A node can update the state \textbf{and} choose the next node in one return value.
        \item Useful when the decision is made inside the node, such as hand-offs between agents.
        \item Replaces most "modify the graph at runtime" ideas: the graph stays fixed, the path is dynamic.
      \end{itemize}
\begin{lstlisting}[language=Python]
from typing import Literal
from langgraph.types import Command

def router(state: RouteState) -> Command[Literal["big_task", "small_task"]]:
    target = "big_task" if state["size"] > 5 else "small_task"
    return Command(update={"path": ["router"]}, goto=target)

builder.add_node("router", router)
builder.add_edge(START, "router")
\end{lstlisting}
\end{frame}
%%%%%%%%%%%%%%%%%%%%%%%%%%%%%%%%%%%%%%%%%%%%%%%%%%%%%%%%%%%%%%%%%%%%%%%%%%%%%%%%
\begin{frame}[fragile]\frametitle{Human-in-the-Loop: Implementation Patterns}
      \begin{itemize}
        \item \textbf{Pattern 1: Breakpoints}:
        \begin{itemize}
            \item Pause execution at specific nodes
            \item Example: Stop before "send\_email" node
        \end{itemize}
        \item \textbf{Pattern 2: Approval Gates}:
        \begin{itemize}
            \item Node requests approval with \texttt{interrupt()}, waits for response
            \item Conditional edge based on approval status
        \end{itemize}
        \item \textbf{Pattern 3: State Editing}:
        \begin{itemize}
            \item Pause, allow human to modify state
            \item Resume with corrected state
        \end{itemize}
        \item \textbf{Pattern 4: Confidence Routing}:
        \begin{itemize}
            \item Low-confidence results go to a human, high-confidence ones continue
            \item Useful for long-running agents
        \end{itemize}
      \end{itemize}
\end{frame}
%%%%%%%%%%%%%%%%%%%%%%%%%%%%%%%%%%%%%%%%%%%%%%%%%%%%%%%%%%%%%%%%%%%%%%%%%%%%%%%%
\begin{frame}[fragile]\frametitle{Human-in-the-Loop: Breakpoint Code}
\begin{lstlisting}[language=Python]
from langgraph.checkpoint.memory import MemorySaver
from langgraph.graph import StateGraph, START, END

# Checkpointing is required for pausing and resuming
memory = MemorySaver()
workflow = StateGraph(EmailState)

workflow.add_node("draft_email", draft_email_node)
workflow.add_node("send_email", send_email_node)
workflow.add_edge(START, "draft_email")
workflow.add_edge("draft_email", "send_email")
workflow.add_edge("send_email", END)

# Compile with a breakpoint BEFORE the send_email node
app = workflow.compile(
    checkpointer=memory,
    interrupt_before=["send_email"])      # pause here for human approval
\end{lstlisting}
\end{frame}
%%%%%%%%%%%%%%%%%%%%%%%%%%%%%%%%%%%%%%%%%%%%%%%%%%%%%%%%%%%%%%%%%%%%%%%%%%%%%%%%
\begin{frame}[fragile]\frametitle{Human-in-the-Loop: Resume and Edit}
\begin{lstlisting}[language=Python]
config = {"configurable": {"thread_id": "1"}}
app.invoke({"email_content": "Meeting request"}, config)

# Execution is paused before send_email
snapshot = app.get_state(config)
print(snapshot.next)              # ('send_email',)
print(snapshot.values["draft"])   # human reviews the draft

# Approve as is: resume from where it stopped
app.invoke(None, config)

# ...or edit the state first, then resume
app.update_state(config, {"draft": "Modified email content"})
app.invoke(None, config)
\end{lstlisting}
\end{frame}
%%%%%%%%%%%%%%%%%%%%%%%%%%%%%%%%%%%%%%%%%%%%%%%%%%%%%%%%%%%%%%%%%%%%%%%%%%%%%%%%
\begin{frame}[fragile]\frametitle{Human-in-the-Loop: Confidence Routing with interrupt()}
      \begin{itemize}
        \item Route low-confidence results to a human node; the graph itself waits, with no \texttt{input()} call.
        \item \texttt{interrupt()} hands a payload to the caller; \texttt{Command(resume=...)} supplies the answer.
      \end{itemize}
\begin{lstlisting}[language=Python]
from langgraph.types import interrupt, Command

def check_confidence(state) -> str:
    return "human_review" if state["confidence"] < 0.8 else "auto_process"

def human_review(state):
    answer = interrupt({"question": "Please confirm or correct",
                        "draft": state["answer"]})
    return {"answer": answer}                 # the human's value

workflow.add_conditional_edges("process_query", check_confidence)

# caller: first invoke pauses and returns the interrupt payload,
# then resume with the human's decision
app.invoke(Command(resume="Corrected answer"), config)
\end{lstlisting}
\end{frame}
%%%%%%%%%%%%%%%%%%%%%%%%%%%%%%%%%%%%%%%%%%%%%%%%%%%%%%%%%%%%%%%%%%%%%%%%%%%%%%%%
\begin{frame}[fragile]\frametitle{Best Practices and Considerations}
      \begin{itemize}
        \item Declare a reducer for every field that several nodes or branches write; the default overwrite loses data.
        \item Design hierarchy and connections of nodes deliberately so the flow is the one you intended.
        \item Use checkpointing, but watch storage overhead and snapshot size for large states.
        \item Use typed state (Pydantic or \texttt{TypedDict} with \texttt{Annotated}) for validation and clarity.
        \item In parallel branches, do not rely on execution order.
        \item Reducers add convenience and some complexity: use them where they pay off.
      \end{itemize}
\end{frame}
%%%%%%%%%%%%%%%%%%%%%%%%%%%%%%%%%%%%%%%%%%%%%%%%%%%%%%%%%%%%%%%%%%%%%%%%%%%%%%%%
\begin{frame}[fragile]\frametitle{Streaming Response in LangGraph}
      \begin{itemize}
        \item Different from traditional LLM word-by-word output
        \item Outputs node data state each time for granular control
        \item Values mode returns complete graph state (total)
        \item Updates mode returns state changes only (incremental)
        \item Messages mode streams LLM tokens as they are produced
        \item Compiled graph is essentially a Runnable component
        \item Enhanced user experience through real-time feedback
      \end{itemize}
\end{frame}
%%%%%%%%%%%%%%%%%%%%%%%%%%%%%%%%%%%%%%%%%%%%%%%%%%%%%%%%%%%%%%%%%%%%%%%%%%%%%%%%
\begin{frame}[fragile]\frametitle{Streaming Modes Implementation}
      \begin{itemize}
        \item Values mode: Complete state after each node
        \item Updates mode: Incremental changes as dictionary
        \item Dictionary keys represent node names
        \item Dictionary values contain state updates
        \item Choose mode based on application requirements
      \end{itemize}
      \begin{lstlisting}[language=Python]
# Values mode: Returns complete state values
inputs = {"messages": [("human", "What are the top 3 results of the 2024 Beijing Half Marathon?")]}

for chunk in agent.stream(inputs, stream_mode="values"):
    chunk["messages"][-1].pretty_print()

# Updates mode: Returns state updates only
for chunk in agent.stream(inputs, stream_mode="updates"):
    print(chunk)
      \end{lstlisting}

\end{frame}
%%%%%%%%%%%%%%%%%%%%%%%%%%%%%%%%%%%%%%%%%%%%%%%%%%%%%%%%%%%%%%%%%%%%%%%%%%%%%%%%
\begin{frame}[fragile]\frametitle{Streaming Example}

      \begin{itemize}
        \item Streaming provides transparency into agent execution
        \item Critical for production applications with users waiting
      \end{itemize}

\begin{lstlisting}[language=Python]
# Instead of invoke(), use stream()
for chunk in app.stream(initial_state, config, stream_mode="updates"):
    # chunk is {node_name: state_update} after each node
    for node_name, update in chunk.items():
        print(f"Node: {node_name}")
        print(f"Update: {update}")

# For token-by-token LLM streaming
async for event in app.astream_events(initial_state, config):
    if event["event"] == "on_chat_model_stream":
        # Stream each token as LLM generates
        print(event["data"]["chunk"].content, end="", flush=True)
      \end{lstlisting}
\end{frame}
